\section{Discussion}\label{sec:discussion}
\subsection{What remains of the index effect}

The replication answers the return question more clearly than the mechanism question. Additions retain a positive average market-adjusted return over 2017--2026, including when influential returns and coincident index events are excluded separately. The later return also exceeds that of the size-screened comparison securities after correcting the overlapping control window. The latest five-year estimate is less precise and more sensitive to individual events. The defensible answer is that a positive addition association remains in the later decade, while its current magnitude is not precisely established. Deletions provide weaker evidence because their intervals permit materially different return magnitudes.

This result is relevant to \citet{greenwood2025} because the disappearance they document is an empirical development, not an automatic consequence of the growth of indexing. Sweden supplies another setting in which to examine the same questions. Reported benchmark-linked purchases and implementation trading are clearly visible, yet the pooled later addition return remains positive. This extends the institutional scope of the question without establishing that otherwise identical Swedish and US stocks respond differently. The periods, selection procedures, investor populations and measurement conventions differ.

The comparison must also distinguish the anchor's pooled S\&P 500 result from its migration analysis. When an addition to one tracked index is a deletion from another, purchases by the destination's trackers need not measure net index-related demand. We have not reconstructed an equivalent historical system of Swedish origin and destination indices. Purchases by other passive funds and the large returns of coincident MSCI events are informative about our sample, but do not prove that offsetting demand is absent elsewhere. The Swedish contribution consequently does not depend on treating every addition as a verified direct addition.

\subsection{What the trading evidence explains}

Fund holdings make the return result economically more informative. On the same later addition events, covered benchmark-linked funds purchase 1.03\% of shares outstanding and CAR averages 2.71\%. Their ratio of 2.63 describes how much return accompanies a percentage point of measured purchases. It is not the slope from an experiment in which demand alone changes. The holdings-report interval also differs from the return window, so the ratio should not be described as an immediate response to orders submitted at implementation.

Two issues prevent that interpretation. First, membership follows firm characteristics and prior performance. News or investor attention can move prices during the window even after market adjustment. Second, observed funds cover only part of the investor population. Omitted institutions can buy or sell, and absent security positions can be inferred as zero. A partial sum of net trades is therefore not necessarily a lower bound on total net index-related demand. Coverage can change the measured ratio even if the response to an otherwise identical order remains unchanged.

The FI comparison strengthens a specific part of the measurement: recent reported purchases have a similar scale through two reporting channels. It does not authenticate every historical zero position or establish constant coverage. The fixed-fund sensitivity also describes a narrower population rather than recovering missing investors. Its mean purchases fall between the two later blocks. Together, these checks support using recent holdings to describe trading, while limiting a claim that rising total passive demand explains a precisely measured historical multiplier decline.

Quoted spreads and daily volume describe the accompanying trading environment. Narrower spreads coexist with more concentrated trading immediately before effectiveness. A small implementation-window return is compatible with investors anticipating the trades or supplying liquidity when they occur. It is also compatible with prices responding earlier to information associated with selection. Closing quotes and daily volume do not distinguish these channels. They show why a small return near implementation need not mean that reconstitution trading has become unimportant.

\subsection{Timing and the nature of the remaining return}

The timing extension changes the description of the historical decline. The first-to-last addition difference is -7.03 percentage points for the full episode and +0.62 points for the common initial seven daily returns. The latter interval, [-1.80, 3.03], establishes neither equality nor an increase. The remaining interval accounts for -7.64 points of the full contrast. The unusually high early total return is principally accumulated after the common initial interval, during episodes that generally allowed more time before implementation. This is a concrete measurement finding even though duration, firm composition and intervening information cannot be separated causally.

The subsequent path addresses a different question. On the fixed later-decade cohort, cumulative CAR is 3.30\% through ED+20. The observed average has not returned to zero by that horizon. This is more informative than treating an insignificant post-event coefficient as proof of no reversal, but it does not identify a permanent index premium. The counterfactual return is unknown, the horizon was examined after the original results, and the latest-block estimate is sensitive to exclusions.

The sister-class comparison asks whether the movement is confined to the security whose status changes. In its selected sample, the index class and sister class have sign-adjusted log market-excess means of 3.18\% and 2.85\%, leaving an imprecise difference of 0.33 points. Shared firm information is one possible explanation. Substitution and arbitrage can also transmit demand between related classes, while voting rights, liquidity and investor clienteles differ. An unchanged index status therefore does not make the sister class unaffected. The small relative estimate neither proves a shared-information explanation nor supplies a causal upper bound on class-specific index demand.

The later-decade path and the separately selected share-class sample describe different aspects of observed returns. Several mechanisms remain compatible with both findings. Documenting returns around changes is a different task from assigning those returns to the orders of particular funds. The former is supported more directly here than the latter.

\subsection{Scope and implications}

Historical notice fields and corporate actions are not fully reconciled. Return availability and the holdings adjustment screen select changing samples, particularly early in the history. Moreover, many securities do not create many independent reviews: the latest block contains nine review clusters, and the research includes several retrospective specifications. The appendix records unresolved source issues and retains unfavorable as well as favorable sensitivities. These limitations are especially consequential for the historical multiplier, less directly for the independently reproduced later return mean.

The estimates are not a tested investment strategy. They neither deduct trading costs nor demonstrate that additions were predictable before the public notice. For research on indexing, the useful finding is that a later addition association coexists with reported benchmark buying and more concentrated trading, while its attribution remains unresolved. The combination of return and holdings evidence extends the anchor question without requiring a structural demand estimate.

A sharper mechanism test would require evidence that separates the competing accounts. Complete historical portfolios would help reconcile omitted positions and changing fund coverage. Class-specific transactions and investor identities would help distinguish direct purchases from substitution across classes. A credible source of membership or demand variation unrelated to firm news would be needed for causal attribution; more controls alone would not provide it. These requirements follow from the findings. The present study establishes what remains in the observed return series and why its timing and investor coverage matter for interpretation.
\FloatBarrier
